\section{Discussion and Limitations}
Softening the date filter does not raise average accuracy on held-out questions. The kernel works as intended: it recovers evidence when the filter finds none and it helps multi-period questions. But it also lets nearby facts crowd out exact ones in single-period, multi-entity questions, and the two effects cancel. The development gain of $+0.077$ vanished on test. The 37-question interval was very wide, and zero-evidence cases were rarer on test. Fixing the settings and the comparisons in advance is what makes this null result trustworthy rather than a tuning artefact. The robustness result is the most encouraging finding, but it rests on one significant test out of six and is within run-to-run variation.

\textbf{Limitations.} (i) With 152 test questions, intervals of about $\pm 0.05$ cannot detect smaller effects. (ii) A single judge model with our own prompt wording produces the main metric. A blind check of 50 judged answers against a human reader is prepared but not yet scored. ROUGE-L and F1 are low for every system ($\approx 0.19$) because gold answers are short. (iii) Each configuration was run once, and the noise test used one noise seed. (iv) One-quarter shifts model only one kind of error; malformed timestamps (9.9\%) get weight 0 under both methods. (v) Six of 24 companies; absolute scores are not comparable with published full-benchmark numbers.
